\section{La salida de Xu Huazhe (许华哲): innovación algorítmica y barreras en la cadena de valor}

Después del cerebro del robot, la entrevista llega a una pregunta delicada pero de gran valor: ¿qué revela la salida de un cofundador? Esta sección no emite juicios sobre rumores; la trata como una ventana para entender cómo se jerarquizan las barreras de entrada de una empresa de robótica.

Cuando se grabó la entrevista, Xu Huazhe, cofundador de Xinghaitu (星海图), estaba por dejar la empresa. El entrevistador insistió: ¿significa esto que, en la etapa actual, la innovación algorítmica no es importante para las empresas de robótica? La respuesta de Gao Jiyang (高继扬) fue clara: la innovación algorítmica es importante, pero no puede existir de forma independiente, separada de toda la cadena de valor de la inteligencia corporizada.

\begin{figure}[H]
\centering
\includegraphics[width=0.96\textwidth]{embodied-ai-value-chain.png}
\caption{Cadena de valor de la inteligencia corporizada y ciclos de difusión: robot completo, clientes, datos, AI infra, algoritmos/modelos.}
\end{figure}

\begin{knowledgebox}{Lectura de la figura: el ciclo de difusión determina las barreras de corto plazo}
Los ciclos de difusión de la figura provienen de las estimaciones hechas en la entrevista: el robot completo y la cadena de suministro requieren entre 12--18 meses; los canales de clientes, al menos 6 meses; el sistema de datos, de 6--12 meses adicionales a partir del robot completo; en cambio, un equipo de primer nivel puede alcanzar un artículo de algoritmos o un método de código abierto en apenas 2--3 meses. Por ello, las barreras de corto plazo no dependen solo del algoritmo, sino de toda la cadena.
\end{knowledgebox}

\subsection{Por qué el algoritmo sigue siendo importante}

Antes se usaron los ciclos de difusión para mostrar que el algoritmo no es necesariamente la barrera más sólida, pero eso no significa que pueda ignorarse. Una lectura más precisa es esta: el algoritmo debe integrarse en activos de ciclo más largo para que la innovación se convierta en barrera.

Gao Jiyang no menospreció el algoritmo. Subrayó que el equipo de algoritmos de Xinghaitu es muy capaz y que seguirá innovando. Pero en el entorno actual, en el que el código abierto y los artículos se difunden muy rápido, el algoritmo por sí mismo es más fácil de replicar; sin el respaldo del robot completo, los datos, la infra, los canales y el valor para el cliente, es difícil que la innovación algorítmica se convierta en una defensa de largo plazo.

\begin{importantbox}{Innovación pragmática}
«Innovación pragmática» no significa dejar de innovar, sino evaluar primero el ROI de la innovación dentro de toda la cadena de valor: cuánto aporta a la estrategia de largo plazo, cuánto aporta a los resultados de corto plazo y si el robot completo, los datos y el ciclo cerrado con los clientes pueden amplificarla.
\end{importantbox}

\subsection{Valores y decisiones organizacionales}

Una vez tratados el algoritmo y las barreras, la pregunta pasa de manera natural a la organización: cuando las personas, la dirección y los recursos ya no coinciden por completo, ¿cómo decide el equipo fundador? Ahí es donde la «innovación pragmática» se pone realmente a prueba.

Gao Jiyang también explicó la salida de Xu Huazhe como un cambio de etapa de la organización: en cada etapa, la empresa necesita distintos tipos de personas para crear valor, y también necesita hacer ajustes con apego a los hechos. Definió los valores como dos tipos de decisiones: ante una elección de dirección, qué se elige y qué no; ante el reparto de beneficios, a quién se le asigna y a quién no. Esto se acerca más a la lógica real de una organización emprendedora que preguntarse «quién tiene la razón y quién no».

\begin{warningbox}{No atribuir la salida de un cofundador a una sola causa}
Los cambios en el equipo fundador pueden deberse a una combinación de dirección, etapa, responsabilidades, valores y la voluntad personal de emprender. Para un observador externo, lo más importante es ver si la organización, tras el ajuste, sigue produciendo resultados, y no fijarse solo en el cambio de personal en sí.
\end{warningbox}

\subsection{Resumen de la sección}

El pasaje sobre la salida de Xu Huazhe, en realidad, deja claras las barreras de una empresa de inteligencia corporizada: el algoritmo es importante, pero debe integrarse en la cadena del robot completo, la cadena de suministro, los datos, la infra, los modelos, la distribución y el valor para el cliente.
